\documentclass[crop=false]{standalone}
% --- ПРЕАМБУЛА ДЛЯ ПОДДЕРЖКИ РУССКОГО ЯЗЫКА И МАТЕМАТИКИ ---
\usepackage[T2A]{fontenc}
\usepackage[utf8]{inputenc}
\usepackage[russian]{babel}
\usepackage{amsmath}
\usepackage{geometry}
\usepackage{amssymb}
\usepackage{graphicx}
\usepackage{float}
\usepackage{import}
\geometry{a4paper, top=2cm, bottom=2cm, left=2cm, right=2cm}
\newcommand{\sidecomment}[1]{\marginpar{\raggedright\scriptsize\itshape #1}}
\newcommand{\iu}{\mathrm{i}}

\newcommand{\Def}{%
    \text{Def}%
    \hspace{-0.85em}% Сдвиг назад ровно на середину (подберите под шрифт)
    \raisebox{0.1ex}{/}% Черта
    \hspace{1em}% Возврат к концу слова + небольшой отступ
}

\renewcommand{\Re}{\operatorname{Re}}
\renewcommand{\Im}{\operatorname{Im}}

\ifstandalone
    \newcommand{\lecturebreak}{} % В отдельном файле лекции разрыв не нужен
\else
    \newcommand{\lecturebreak}{\clearpage} % В полной книге - начинаем с новой страницы
\fi

\newcounter{lecturenumber} % Создаем счетчик для нумерации лекций
\newcommand{\lecture}[2]{
    \lecturebreak
    \refstepcounter{lecturenumber} % Увеличиваем счетчик на 1
    \par\vspace{2em}\noindent % Отступ сверху и убираем абзацный отступ
    {\Large\bfseries % Делаем заголовок крупным и жирным
    Лекция \thelecturenumber: #1 % Выводим "Лекция N: Тема"
    \hfill % Растягиваем пробел, чтобы дата была справа
    \large\textit{#2}} % Выводим дату курсивом
    \par\vspace{1em} % Небольшой отступ снизу
    \addcontentsline{toc}{section}{Лекция \thelecturenumber: #1} % Добавляем в оглавление
    \par
}

\newcommand{\TwoCols}[2]{%
  \par\noindent % Начать с новой строки без отступа
  \begin{minipage}[t]{0.48\textwidth}
    #1
  \end{minipage}% <-- Важный '%' для удаления пробела
  \hfill % Растягивающийся пробел между колонками
  \begin{minipage}[t]{0.48\textwidth}
    #2
  \end{minipage}% <-- Важный '%' для удаления пробела
  \par%
}

\pagestyle{plain} % Включаем нумерацию страниц\
\begin{document}
\section{Уравнения Лагранжа второго рода}
\keypoint уравнения Лагранжа первого рода это способ решать дифуры с ограничениями, в механике обычно не используются

построение уравнений Лагранжа первого рода
\[
\begin{cases}
\sum_{k=1}^N (\bar F - m_k \bar w_k) \delta \bar r_k =0 \\ f_s(\bar r_k) = 0\\ f_p(\bar r_k) = 0
\end{cases}
\]

\noindent \keypoint второй род

для голономных систем:
\[0 = \sum_{k=1}^{N} (\bar F_k - m_k \bar w_k) \delta \bar r_k = \sum_{k=1}^{N} (\bar F_k - m_k \bar w_k) \sum_{i=1}^{n} \frac{\partial \bar r_k}{\partial q_i} \delta q_i \]
\[\sum_{i=1}^{n} \left( \underbracket{\sum_{k=1}^{N} \bar F_k \frac{\partial \bar r_k}{\partial q_i}}_{Q_i} -  \sum_{k=1}^{N} m_k \bar w_k \frac{\partial \bar r_k}{\partial q_i} \right) \delta q_i = 0\]
$\delta q_i$ независимы, то есть каждая скобка равна нулю, то есть
\[\sum_{k=1}^N m_k \bar w_k \frac{\partial \bar r_k}{\partial q_i} = Q_i\]
\[
\begin{aligned}
Q_i = \sum_{k=1}^N m_k \bar w_k \frac{\partial \bar r_k}{\partial q_i} &= \sum_{k=1}^{N} m_k \frac{d \bar v_k}{d t} \frac{\partial \bar r_k}{\partial q_i} = \sum_{k=1}^{N} m_k  \left( \frac{d}{d t} \left(\bar v_k, \frac{\partial \bar r_k}{\partial q_i}\right) - \bar v_k \frac{d}{d t} \frac{\partial \bar r_k}{\partial q_i} \right) \\
&= \frac{d}{d t} \sum_{k=1}^{N} \frac{m_k}{2} \frac{\partial v_k^2}{\partial \dot q_i} - \sum_{k=1}^{N} \frac{m_k}{2} \frac{\partial v_k^2}{\partial q_i} \\ &= \frac{d}{d t} \frac{\partial}{\partial \dot q_i}\sum_{k=1}^{N} \frac{m_k v_k^2}{2} - \frac{\partial}{\partial q_i} \sum_{k=1}^{N} \frac{m_k v_k^2}{2} = \frac{d}{d t} \frac{\partial}{\partial \dot q_i} T - \frac{\partial}{\partial q_i} T
\end{aligned}
\]

\[\frac{d}{d t} \frac{\partial T}{\partial \dot q_i} - \frac{\partial T}{\partial q_i} = - \frac{\partial \Pi}{\partial q_i} \qquad \frac{d}{d t} \frac{\partial T}{\partial \dot q_i} - \frac{\partial (T - \Pi)}{\partial q_i} =0 \qquad \frac{\partial \Pi}{\partial \dot q_i} = 0\]
\[\frac{d}{d t} \frac{\partial(T -\Pi)}{\partial \dot q_i} - \frac{\partial (T - \Pi)}{\partial q_i} =0 \qquad L := T - \Pi\]
\[\boxed{\frac{d}{d t} \frac{\partial L}{\partial \dot q_i} - \frac{\partial L}{\partial q_i} =0}\]
\Def $L$ — Лагранжиан

Это и есть уравнение Лагранжа второго рода

Это не уравнения в частных производных, их $n$ штук, разрешимость уравнений Лагранжа следует из основной теоремы Лагранжевого формализма

\Def Координатное пространство — $q_1, q_2, \dots q_n$, точка характеризует положение системы

\Def Фазовое пространство — $q_1 , q_2, \dots, q_n , \dot q_1, \dots, \dot q_n$, точка характеризует состояние системы

В современной физике уравнения Лагранжа выводятся из принципа наименьшего действия, из него же можно вывести второй закон Ньютона, и общую теорию относительности

\Def Натуральные системы — системы, в которых Лагранжиан является разность кинетической и потенциальной энергии
\[L = T - \Pi\]

\keypoint Пример

Математический маятник: длина нити $l$, масса груза $m$, угол от вертикал $\varphi$
$x$ вниз, $y$ вправо, центр в подвесе, на груз действует сила $\bar F = \begin{pmatrix}
    mg \\ 0
\end{pmatrix}$, $\Pi = - mgx$
\[x = l \cos \varphi \quad y = l \sin \varphi \implies \Pi = - m g l \cos \varphi\]
\[T = \frac{m v^2}{2} \qquad \bar r = \begin{pmatrix}
    x \\ y
\end{pmatrix} \implies \bar v = \begin{pmatrix}
    \dot x \\ \dot y
\end{pmatrix} \quad T = \frac{m}{2} \left(\dot x^2 + \dot y^2\right) = \frac{m}{2} l^2 \dot \varphi^2\]
\[L = T - \Pi = \frac{m l^2 \dot \varphi^2}{2} + m g l \cos \varphi\]
\[\frac{d}{d t} \frac{\partial}{\partial \dot \varphi} \left(\frac{m l^2 \dot \varphi^2}{2} + m g l \cos \varphi\right) - \frac{\partial }{\partial \varphi} \left( \frac{m l^2 \dot \varphi^2}{2} + m g l \cos \varphi \right) = 0\]
\[\frac{d}{d t} (m l^2 \dot \varphi) + m g l \sin \varphi = 0\]
\[m l^2 \ddot \varphi + m g l \sin \varphi = 0 \qquad \ddot \varphi + \frac{g}{l} \sin \varphi = 0\]


\end{document}